\documentclass[14pt,a4paper]{extarticle}

\usepackage[T2A]{fontenc}
\usepackage[utf8]{inputenc}
\usepackage[russian]{babel}
\usepackage[a4paper,margin=2cm]{geometry}
\usepackage{amsmath,amssymb}
\usepackage{graphicx}
\usepackage{booktabs}
\usepackage{float}
\usepackage{array}
\usepackage{caption}
\usepackage{hyperref}
\usepackage{setspace}

\onehalfspacing
\emergencystretch=3em
\graphicspath{{../outputs/lab5/plots/}}

\hypersetup{
    colorlinks=true,
    linkcolor=black,
    urlcolor=blue
}

\newcommand{\plotplaceholder}[2]{
    \begin{figure}[H]
        \centering
        \IfFileExists{../outputs/lab5/plots/#1}{
            \includegraphics[width=0.88\textwidth]{#1}
        }{
            \fbox{\parbox[c][5cm][c]{0.82\textwidth}{\centering Файл \texttt{\detokenize{#1}} создается запуском\\\texttt{python lab5/run.py}}}
        }
        \caption{#2}
    \end{figure}
}

\newcommand{\fitreporttable}[1]{
    \resizebox{\textwidth}{!}{#1}
}

\title{Лабораторная работа №5\\Оптимизация в задаче регрессии}
\author{
    Новиков Алексей Георгиевич, группа М3233\\
    Душкин Александр Алексеевич, группа М3233
}
\date{2026}

\begin{document}

\begin{titlepage}
    \begin{center}
        {\large Отчет по лабораторной работе №5}\\[0.5cm]
        {\LARGE \textbf{Полиномиальная регрессия и методы оптимизации}}\\[2cm]

        \begin{flushright}
            Выполнили:\\[0.2cm]
            Новиков Алексей Георгиевич, группа М3233\\
            Душкин Александр Алексеевич, группа М3233\\[0.5cm]
        \end{flushright}

        \vfill
        Санкт-Петербург\\[0.2cm]
        2026
    \end{center}
\end{titlepage}

\tableofcontents
\newpage

\section{Постановка задачи}

Цель работы --- рассмотреть задачу одномерной регрессии как задачу оптимизации и сравнить несколько способов минимизации функции потерь. В работе использовались полиномиальные модели степеней от 1 до 5, два искусственных набора данных, варианты без регуляризации, с L1, L2 и Elastic Net, а также несколько оптимизаторов: аналитическое решение для линейной модели, SGD, mini-batch GD, Гаусс--Ньютон и Левенберг--Марквардт.

Для модели степени $d$ использовалась матрица признаков
\[
    X = [1,z,z^2,\ldots,z^d],
\]
где $z$ --- нормированное значение исходного признака. Базовая функция потерь:
\[
    L(w)=\frac1m\|Xw-y\|^2.
\]
При регуляризации добавлялись члены
\[
    \lambda_1\sum_{i=1}^{d} \sqrt{w_i^2+\delta^2}
    \quad\text{и}\quad
    \lambda_2\sum_{i=1}^{d} w_i^2.
\]
Свободный коэффициент $w_0$ в этих экспериментах не регуляризировался: штраф применялся только к коэффициентам при степенях нормированного признака. Для L1 использовалась сглаженная аппроксимация модуля с малым $\delta$, чтобы сохранять дифференцируемость функции потерь для градиентных и гаусс--ньютоновских методов.

\textbf{Промежуточный вывод.} Регрессия удобна тем, что одна и та же задача допускает разные методы. Для линейного случая есть аналитическая формула, для стохастических методов важны батчи и шаг, а методы Гаусса--Ньютона используют структуру суммы квадратов через якобиан остатков.

\section{Данные и модели}

Первый набор данных почти линейный:
\[
    y=1.6x-0.7+0.1\sin(3x)+\varepsilon,\qquad \operatorname{Var}(\varepsilon)=0.04.
\]
Второй набор данных существенно нелинейный:
\[
    y=0.35x^3-0.8x+\sin(2.8x)+1.5\exp\left(-\frac12\left(\frac{x-1}{0.35}\right)^2\right)+\varepsilon,
\]
где дисперсия шума равна $0.09$.

\textbf{Промежуточный вывод.} Почти линейная зависимость нужна как контрольный пример: даже модель первой степени уже близка к данным. Нелинейный набор показывает, зачем увеличивать степень полинома и зачем нужна регуляризация, когда модель становится гибкой.

\section{Методы оптимизации}

\subsection{Аналитическое решение}

Для линейной модели без регуляризации используется решение через выборочные средние. Такой метод применим только к степени 1 и без L1/L2 членов.

\textbf{Промежуточный вывод.} Аналитическое решение почти бесплатно по числу итераций, но оно узкоспециализировано. Как только меняется степень модели или добавляется регуляризация, этот путь уже не подходит.

\subsection{SGD и mini-batch GD}

Стохастический градиентный спуск обновляет веса по одному объекту, а mini-batch GD использует подвыборку:
\[
    w_{t+1}=w_t-\alpha_t\nabla L_{B_t}(w_t).
\]
В экспериментах использовалось убывание шага по правилу $\alpha_t=\alpha_0/\sqrt{t+1}$.

\textbf{Промежуточный вывод.} Малый batch дает шумный, но частый сигнал о направлении. Большой batch ближе к полному градиенту, зато одно обновление становится менее локальным по данным. Поэтому сравнение batch size лучше смотреть не только по эпохам, но и по числу градиентных вычислений.

\subsection{Гаусс--Ньютон и Левенберг--Марквардт}

Для суммы квадратов остатков можно использовать якобиан $J$:
\[
    (J^TJ)p=-J^Tr.
\]
Метод Левенберга--Марквардта добавляет демпфирование:
\[
    (J^TJ+\mu I)p=-J^Tr.
\]

\textbf{Промежуточный вывод.} Гаусс--Ньютон быстро работает, когда квадратичная аппроксимация остатков хороша. Демпфирование в Левенберге--Марквардте делает шаг осторожнее: при плохом локальном приближении метод ведет себя ближе к градиентному, а при хорошей модели --- ближе к Гауссу--Ньютону.

\section{Результаты}

\subsection{Сравнение методов}

\begin{table}[H]
\centering
\caption{Сводка по основным оптимизаторам}
\fitreporttable{
\begin{tabular}{lrrrrr}
\toprule
Метод & Успешных запусков & Лучшее $f$ & Средн. итераций & $N_f$ & $N_g$\\
\midrule
AnalyticalLinearRegression1D & 2 & 0.0394 & 1.00 & 1.00 & 1.00\\
GaussNewton & 18 & 0.0385 & 23.00 & 24.00 & 24.00\\
LevenbergMarquardt & 22 & 0.0385 & 13.91 & 14.91 & 11.05\\
MiniBatchGradientDescent & 0 & 0.0618 & 80.00 & 81.00 & 1063.07\\
StochasticGradientDescent & 0 & 0.0394 & 80.00 & 81.00 & 9681.00\\
\bottomrule
\end{tabular}
}
\end{table}

\textbf{Промежуточный вывод.} Методы Гаусса--Ньютона и Левенберга--Марквардта быстрее приходят к малой ошибке, потому что используют якобиан всей модели. SGD и mini-batch не достигали строгой точности за 80 эпох, но на части задач давали близкие значения функции потерь.

\subsection{Влияние степени полинома}

\begin{table}[H]
\centering
\caption{Лучшие значения функции потерь для Levenberg--Marquardt}
\fitreporttable{
\begin{tabular}{lrrrrr}
\toprule
Данные & degree 1 & degree 2 & degree 3 & degree 4 & degree 5\\
\midrule
near\_linear & 0.03942 & 0.03936 & 0.03928 & 0.03861 & 0.03853\\
nonlinear & 1.17604 & 1.15854 & 0.82320 & 0.80880 & 0.28341\\
\bottomrule
\end{tabular}
}
\end{table}

\textbf{Промежуточный вывод.} Для почти линейных данных увеличение степени дает небольшой выигрыш: основная структура уже описана первой степенью. Для нелинейной зависимости степень 5 резко улучшает результат, потому что модель наконец получает достаточно гибкости для пика и осцилляции.

\subsection{Регуляризация}

\begin{table}[H]
\centering
\caption{Регуляризация для nonlinear degree 5, Levenberg--Marquardt}
\fitreporttable{
\begin{tabular}{lrr}
\toprule
Регуляризация & $f$ & Итерации\\
\midrule
нет & 0.28341 & 4\\
L1, $\lambda=0.001$ & 0.29050 & 4\\
L1, $\lambda=0.01$ & 0.35078 & 5\\
L1, $\lambda=0.1$ & 0.63649 & 53\\
L1, $\lambda=1$ & 0.88710 & 58\\
L2, $\lambda=0.001$ & 0.29969 & 4\\
L2, $\lambda=0.01$ & 0.39649 & 4\\
L2, $\lambda=0.1$ & 0.57235 & 3\\
L2, $\lambda=1$ & 0.69107 & 3\\
Elastic Net, $\lambda=0.01$ & 0.44208 & 4\\
Elastic Net, $\lambda=1$ & 0.93095 & 85\\
\bottomrule
\end{tabular}
}
\end{table}

\textbf{Промежуточный вывод.} Регуляризация увеличивает оптимизируемую функцию, потому что штраф входит прямо в значение $f$. Это не ошибка: смысл штрафа --- не уменьшить training loss любой ценой, а сдержать коэффициенты и сделать модель менее резкой.

\section{Графики}

\plotplaceholder{PolynomialRegressionObjective__dataset_kind_near_linear_degree_5_n_points_120_noise_variance_0.04_seed_11_x_range_-3.0_3.0_fit.png}{Аппроксимация почти линейных данных полиномом степени 5}
\textit{Комментарий к графику.} Истинная зависимость почти линейна, поэтому высокая степень нужна мало. Если кривая почти совпадает с низкостепенной моделью, это говорит, что дополнительные коэффициенты в основном подстраивают мелкие колебания и шум.

\plotplaceholder{PolynomialRegressionObjective__dataset_kind_nonlinear_degree_5_n_points_120_noise_variance_0.09_seed_17_x_range_-2.5_2.5_fit.png}{Аппроксимация нелинейных данных полиномом степени 5}
\textit{Комментарий к графику.} Нелинейный набор содержит пик и осцилляцию. Полином высокой степени лучше повторяет форму данных, но именно здесь возрастает риск лишней гибкости и становится полезной регуляризация.

\plotplaceholder{PolynomialRegressionObjective__dataset_kind_nonlinear_degree_5_n_points_120_noise_variance_0.09_seed_17_x_range_-2.5_2.5_history_loss.png}{История loss для nonlinear degree 5}
\textit{Комментарий к графику.} Методы на всей выборке обычно дают более гладкую кривую loss, потому что используют полный градиент или якобиан. Стохастические методы движутся шумнее: каждый шаг видит только часть данных.

\plotplaceholder{PolynomialRegressionObjective__dataset_kind_nonlinear_degree_5_n_points_120_noise_variance_0.09_seed_17_x_range_-2.5_2.5_history_loss_by_grad_calls.png}{Loss в зависимости от числа вычислений градиента}
\textit{Комментарий к графику.} Такая ось честнее показывает цену стохастических методов. SGD делает много дешевых обновлений, но по счетчику градиентов видно, что путь к малой ошибке может быть дорогим.

\plotplaceholder{lab5_batch_size_comparison_by_epoch.png}{Влияние размера batch по эпохам}
\textit{Комментарий к графику.} Малые батчи обновляют веса чаще и шумнее, крупные батчи ближе к полному градиенту. Поэтому график по эпохам показывает не только качество направления, но и то, сколько обновлений было сделано внутри эпохи.

\plotplaceholder{lab5_regularization_empirical_risk.png}{Эмпирический риск при разных регуляризациях}
\textit{Комментарий к графику.} При росте регуляризации модель жертвует точностью на обучающей выборке ради более сдержанных коэффициентов. Если риск растет слишком сильно, штраф уже подавляет полезную гибкость модели.

\plotplaceholder{lab5_regularized_coefficients.png}{Коэффициенты при регуляризации}
\textit{Комментарий к графику.} L1 стремится занулять менее важные коэффициенты, L2 уменьшает их более плавно, а Elastic Net сочетает оба эффекта. По столбцам коэффициентов видно, как штраф меняет саму модель, а не только итоговый loss.

\plotplaceholder{lab5_method_comparison_loss.png}{Сравнение методов по итоговой функции потерь}
\textit{Комментарий к графику.} Методы Гаусса--Ньютона и Левенберга--Марквардта выигрывают там, где структура суммы квадратов хорошо описывает задачу. Градиентные методы проще, но им нужны дополнительные эпохи и аккуратный шаг.

\plotplaceholder{lab5_method_comparison_elapsed_seconds.png}{Сравнение методов по среднему времени работы}
\textit{Комментарий к графику.} Время работы дополняет счетчики итераций и градиентов: аналитическое решение почти мгновенно, методы Гаусса--Ньютона тратят больше линейной алгебры на шаг, а стохастические методы делают много дешевых обновлений. Поэтому метод с малым числом эпох не обязательно оказывается самым дешевым по фактическому времени.

\section{Полная таблица экспериментов}

Полная сводка 264 запусков хранится в \path{../outputs/lab5/tables/summary.csv}. Ниже приведена печатная версия таблицы, разбитая на масштабируемые блоки. Дополнительные CSV с историями, коэффициентами и данными сохранены в той же папке.

\input{../outputs/lab5/tables/summary_report_scaled.tex}

\section{Общий вывод}

В лабораторной работе регрессия рассматривалась именно как задача оптимизации. Для линейной модели без регуляризации аналитическое решение оказалось самым дешевым, но оно применимо только в узком случае. Методы Гаусса--Ньютона и Левенберга--Марквардта лучше использовали структуру суммы квадратов и быстро сходились на большинстве полиномиальных моделей.

SGD и mini-batch GD не достигали строгого критерия за заданное число эпох, но показали важную практическую особенность: качество зависит от правила шага и размера батча. На нелинейных данных увеличение степени полинома заметно улучшало аппроксимацию, а регуляризация управляла компромиссом между точностью на обучении и величиной коэффициентов.

\end{document}
